\documentclass[../../../include/open-logic-section]{subfiles}

\begin{document}

\olfileid[fr]{his}{set}{infinitesimals}
\olsection{Infinitésimaux et dérivation}

Newton et Leibniz découvrirent indépendamment le calcul
infinitésimal à la fin du XVIIe siècle. Une application
particulièrement importante en était la \emph{dérivation}. En
première approximation, celle-ci cherche à définir le «~taux de
variation~», ou la pente, d'une fonction en un point.

Voici une illustration concrète de cette idée. Considérons la
fonction $f(x) = \nicefrac{x^2}{4} + \nicefrac{1}{2}$, représentée
ci-dessous en noir :
\begin{center}
	\begin{tikzpicture}[scale=1]
	\draw[->, gray] (-1,0) -- (4.25,0) node[right] {$x$};
	\draw[->, gray] (0,-0.25) -- (0,5.25) node[above] {$f(x)$};

	\foreach \x/\xtext in {1/1, 2/2, 3/3, 4/4}
	\draw[shift={(\x,0)}, gray] (0pt,2pt) -- (0pt,-2pt) node[below] {$\xtext$};

	\foreach \y/\ytext in {1/1, 2/2, 3/3, 4/4, 5/5}
	\draw[shift={(0,\y)}, gray] (2pt,0pt) -- (-2pt,0pt) node[left] {$\ytext$};

	\draw[oldiagcolorC] (0.5, 9/16) -- (3.5, 57/16) -- (3.5, 9/16)--cycle;
	\draw[oldiagcolorD] (.5, 9/16) -- (2.5, 33/16) -- (2.5, 9/16) -- cycle;
	\draw[oldiagcolorE] (.5, 9/16) -- (1.5, 17/16) -- (1.5, 9/16) -- cycle;
	\draw[thick] (-1,.75) parabola bend (0,0.5) (4,4.5);
	\end{tikzpicture}
\end{center}
Supposons que nous voulions trouver la pente de la courbe au point
$c = \nicefrac{1}{2}$. Nous traçons d'abord un triangle dont
l'hypoténuse donne une approximation de cette pente, par exemple
le triangle rouge de la figure. Si la base de ce triangle a
pour longueur $\beta$, sa hauteur vaut
$f(\nicefrac{1}{2}+\beta) - f(\nicefrac{1}{2})$. La pente de
l'hypoténuse est donc :
\[
\frac{f(\nicefrac{1}{2}+\beta) - f(\nicefrac{1}{2})}{\beta}.
\]
La pente de l'hypoténuse du triangle !!{colorC}, dont la base
mesure~$3$, vaut ainsi exactement~$1$. L'hypoténuse d'un triangle
plus petit, le triangle !!{colorD} de base~$2$, fournit une meilleure
approximation : sa pente est $\nicefrac{3}{4}$. Un triangle encore
plus petit, le triangle !!{colorE} de base~$1$, donne une
approximation meilleure encore : sa pente est $\nicefrac{1}{2}$.

Des triangles de plus en plus petits fournissent des approximations
de plus en plus précises. Nous pourrions donc dire ceci :
l'hypoténuse d'un triangle dont la base a une longueur
\emph{infinitésimale} donne la pente de la courbe au point
$c = \nicefrac{1}{2}$ lui-même. Nous obtiendrions ainsi une formule
pour la dérivée (première) de $f$ au point $c$ :
\[
{f'}(c) = \frac{f(c+\beta) - f(c)}{\beta}
\text{ où $\beta$ est infinitésimal.}
\]
À peu de chose près, c'est ce que disaient Newton et Leibniz.

Mais, dès lors, nous devons leur demander : qu'est-ce qu'un
\emph{infinitésimal} ? Un sérieux dilemme apparaît. Si $\beta = 0$,
alors $f'$ n'est pas définie, puisqu'il faudrait diviser par~$0$.
Mais si $\beta > 0$, nous n'obtenons qu'une
\emph{approximation} de la pente, et non la pente elle-même.

Cette objection n'est pas anachronique. Voici comment Berkeley
critiquait les disciples de Newton :
\begin{quote}
J'admets que l'on puisse faire désigner aux signes quelque chose ou
rien du tout, et donc que, dans la notation initiale $c + \beta$,
$\beta$ ait pu signifier soit un accroissement, soit rien. Mais,
quel que soit le sens que vous lui donniez, vous devez raisonner
en accord avec ce sens et ne pas jouer sur une double signification :
ce serait là un sophisme manifeste.
(\citealt[\S{}XIII]{Berkeley1734}, variables adaptées à la notation
précédente)
\end{quote}
Pour défendre le calcul infinitésimal contre Berkeley, on pourrait
répondre que parler d'«~infinitésimaux~» n'est qu'une image. On
pourrait dire qu'un triangle \emph{vraiment petit} donne une
approximation \emph{suffisamment bonne} de la tangente. Berkeley
avait également une réponse à cela : cette approximation peut
suffire en ingénierie, mais elle compromet le \emph{statut} des
mathématiques, car
\begin{quote}
on nous dit que \emph{in rebus mathematicis errores qu\`{a}m minimi
non sunt contemnendi}. [En mathématiques, même les plus petites
erreurs ne doivent pas être négligées.] \citep[\S{}IX]{Berkeley1734}
\end{quote}
Le passage en italique est une citation presque littérale de la
\emph{Quadrature des courbes} (1704) de Newton.

Les objections philosophiques de Berkeley sont d'une grande acuité.
Le calcul infinitésimal n'en fut pas moins une entreprise
extraordinairement féconde, et les mathématiciens continuèrent à
l'utiliser sans tomber dans l'erreur.

\end{document}
